\section{Reed--Solomon Projection And MCA}\label{sec:rs-mca-projection}

The concrete projection chain belongs here because its endpoint is the
Reed--Solomon received-word surface and the mutual-correlated-agreement theorem.
This section uses Grassmannian data in a different role from the HCP
Grassmannian of Sections~\ref{sec:hcp-definition}--\ref{sec:mathematical-framework}.
Earlier, \(W\in\Gr(k,R^n)\) parameterizes hidden harmonic kernels and
projection Laplacians \(\Delta_W=I-P_W\). Here, after the hidden/geometric
object has already been accepted, a frame for the Grassmannian point is used as
the source of a lossy proximity projection through Plucker coordinates and
actual cochain commitments toward a finite Reed--Solomon target.

\subsection{Projection Chain}

The staged chain has the following shape:
\[
\begin{aligned}
  &\text{FrameWitness} \\
  &\quad\to \text{Projective Plucker shadow} \\
  &\quad\to \text{committed cochain/local syndrome} \\
  &\quad\to \text{plane shadow} \\
  &\quad\to \text{finite-field received word} \\
  &\quad\to \text{canonical Reed--Solomon target}.
\end{aligned}
\]
Each stage has local proximity gates and certificates. The orchestrating
projection-chain theorem composes these local certificates without rechecking
them. These are proximity-specific projection gates, not the main \HE{}
acceptance layer from Section~\ref{sec:mathematical-framework}. Their role is
to certify the lossy path into the Reed--Solomon endpoint after the main
hidden/geometric object has already been accepted.
Thus the Plucker shadow is not part of the definition of HCP and is not the
source of hidden-operator hardness. It is the first public projective encoding
used by the proximity proof.

The semantic roles in this chain are separated. The private witness consists of
the fibration data, frame choices, paths, and openings. The shadow is the
Plucker/cochain/proximity data derived from that witness. A hash or root is only
a compressed public handle to canonical bytes for the shadow, not the
mathematical shadow itself. A certificate or SNARK proves that the private
witness produced the accepted shadow and satisfied the stage invariants.

\subsection{Projection-Stage Gates}

The source and Plucker stages carry metric and public-shadow invariants.
We use the usual Grassmannian and Plucker-coordinate algebra: a frame
\((w_1,\ldots,w_r)\) determines the exterior coordinate
\([w_1\wedge\cdots\wedge w_r]\in\mathbb P(\bigwedge^r V)\), with the
quadratic Plucker relations as the projective Grassmannian equations.
The source disagreement score \(\rho(W)\) is required to be nonnegative and
invariant under the public Plucker quotient. The Plucker metric is defined only
on the admissible projective chart, and its distance must be nonnegative and
frame-quotient invariant. Public-shadow respect says that equality of the
public Plucker shadow determines the same committed cochain, plane shadow, and
received word along this projection path.

The Plucker-to-cochain stage carries the local-syndrome interface and consumes
the actual cochain commitment layer of Section~5. Its target is not an abstract
stand-in for a cochain; it is the existing committed payload
\[
  (q,n,c,\tau)
\]
with supported modulus, declared edge count, edge values, and topology hash.
A concrete encoder may send valid Plucker frame data to such a well-formed
payload with the expected edge count, but the storage-native statement is
relational: an observed or proof-carrying cochain commitment \(c\) realizes the
public Plucker shadow. This is the boundary formalized by the public-binding
surface: the chain/storage layer accepts the actual commitment payload, while
the proximity proof supplies the statement that this payload realizes the
Plucker shadow.

Syndrome adequacy then states that a well-formed committed cochain \(c\) with
vanishing local syndrome \(\sigma(c)=0\) is a global-section candidate for the
chosen sheaf. In the concrete bridge, the local syndrome is the
commitment-sum witness; global candidates have zero commitment sum. The point
is that the proximity proof passes through real cochain commitments and their
well-formedness/topology-binding interface, not through a separate synthetic
cochain carrier.

The cochain-to-plane and plane-to-RS stages carry semantic-alignment and
non-collapse conditions. Semantic alignment requires valid frames, the plane
shadow, and the received word to agree with the extraction maps. Plane
extraction records the cochain well-formedness, expected edge count, and
selected-index bounds needed to form the two-coordinate shadow. The far-region
conditions assert calibrated non-collapse from the Plucker source region,
through the committed cochain, to the final Reed--Solomon received word.

\subsection{Canonical Reed--Solomon Endpoint}

The proximity endpoint is stated over the paper-style alphabet shape
\(C \subseteq (\F^s)^n\). The canonical \HE{} evaluation is the ordinary
Reed--Solomon slice at width \(s = 1\), using the standard polynomial
evaluation-code convention.
The canonical endpoint fixes
\[
  \F=\mathbb Z/12289\mathbb Z,\qquad n=4096,\qquad k=1024,
\]
so the rate is \(1/4\). The evaluation domain is a certified smooth power
domain: the chosen generator has order \(12288\), and cubing it gives an
element of order \(4096\). Hence the evaluation points are
\[
  1,\alpha,\alpha^2,\ldots,\alpha^{4095},
\]
with \(\alpha^{4096}=1\) and no earlier repetition. The canonical radius used
for the below-Johnson comparison is the absolute Hamming radius
\[
  2047.
\]
The required finite checks are the field arithmetic, domain cardinality,
smoothness, evaluation-point compatibility, and absolute-radius normalization
for this endpoint; the deep list-decoding and MCA estimates enter only as cited
code-theoretic theorem slots.

\subsection{Mutual Correlated Agreement}

\begin{definition}[MCA bad-challenge fraction]\label{def:mca-bad-challenge}
  Let \(C \subseteq (\F^s)^n\) be a block code and let
  \(\ell(\gamma)=f_1+\gamma f_2\) be an affine line of received words. For a
  radius \(\delta\), the bad challenges are those \(\gamma \in \F\) for which
  \(\ell(\gamma)\) is \(\delta\)-close to \(C\) on a large coordinate set, but
  the pair \((f_1,f_2)\) is not correspondingly close to the two-fold
  interleaved code \(C^{\oplus 2}\) on that same set. The MCA error is the
  fraction of bad challenges in \(\F\).
\end{definition}

For the canonical scalar Reed--Solomon slice \(s=1\), ordinary scalar RS words
are identified with words in \((\F^1)^n\), and the bad-challenge predicate
agrees with the paper-alphabet predicate.

\begin{theorem}[Canonical RS MCA bound]\label{thm:canonical-rs-mca}
  For the canonical \HE{} Reed--Solomon endpoint at the below-Johnson radius,
  every projection-induced final affine line has bad-challenge fraction at most
  \[
    \frac{\tau}{|\F|},
  \]
  where \(\tau\) is the cited MCA threshold for the canonical endpoint.
\end{theorem}

The deep list-decoding and MCA estimates are external theorem inputs: the
Johnson list-decoding bound, the count and fractional MCA bounds, and the paper
Definition~4.3 width-one MCA bound. The internal finite part is the endpoint
plumbing, the \(s=1\) equivalence, and the specialization to projection-induced
lines.
The project-specific list-decoding and MCA input is the cited Proximity Prize
theorem boundary \cite{proximityPrize2026}.

\subsection{Corollary}

\begin{corollary}[Verifier-facing RS/MCA projection]\label{cor:verifier-rs-mca}
  Given an accepted hidden/geometric witness, a certified
  Plucker-to-cochain-to-plane-to-RS projection chain, the required hard
  projection-gate evidence, and a canonical Johnson-radius far-separation
  witness, the public verifier obtains a finite Reed--Solomon received word
  whose canonical MCA bad-challenge fraction is bounded by the cited threshold
  \(\tau/|\F|\).
\end{corollary}

This corollary is a proximity statement. It says that accepted hidden/geometric
state can be projected to a finite public RS/MCA target with calibrated
non-collapse at the selected far region. It does not say that the projection is
injective, and it does not turn the RS/MCA endpoint into the source of HCP
hardness.

\subsection*{Lean Correspondence}

The proximity and finite-code endpoint in this section corresponds to the
following Lean declarations:
\lean{Proximity.GrassmannianPlucker}, \lean{Proximity.CochainBridge},
\lean{Proximity.PluckerCochain}, \lean{Proximity.CochainPlane},
\lean{Proximity.ProjectionChain}, \lean{ProjectionChain.Chain.toEvaluationProjection},
\lean{Proximity.Gates}, \lean{Proximity.GateEvidence},
\lean{CanonicalRS.BridgeData}, \lean{CanonicalRS.CitedBounds},
\lean{CanonicalRS.MCA}, and
\lean{CanonicalRS.MCA.projectionBadChallengeFraction_le}.
